\chapter*{Introduction: Is there a wall here?}\label{ch:intro}
\addcontentsline{toc}{chapter}{Introduction: Is there a wall here?}\markboth{Introduction}{}
A pulse is emitted to the left of a plane and detected to its right. Draw a line on the spacetime diagram at the plane: nothing happens to the pulse. Now give the left and right halves of the diagram to two different physicists. Ask each for a retarded solution with the field held at zero on the line. Each calculation is legitimate, but the detector on the right reads zero. A line in a drawing has become a reflecting wall in an initial-boundary value problem. How can the two physicists use their answers to remove the wall without looking up the solution on the uncut spacetime?

They need one more experiment that each can perform independently. With no bulk source on their own side, prescribe a small history of the field at the cut and record its outward response. The two response functions say how to vary that common history until no uncompensated force acts at the plane. The value required on the plane at noon must not depend on what the source will do tomorrow. This last condition makes the experiment Lorentzian. We will calculate this example, including the first arrival at the detector, in Chapter~\ref{ch:08}.

\section*{Two experiments that run in opposite directions}
In a \emph{cutting test} the whole theory is known. Restrict its fields, equations and observables to chosen pieces; sew them back and compare. In a \emph{construction} the pieces are independently given. One must first find out whether their faces, bundles and fixed backgrounds can be identified. There might be no possible whole, or several different wholes. In particular, two local descriptions in different gauges need not have literally equal values on their common face. They need an identification of their fibers. Holding such an identification fixed is a choice of geometric problem, not the insertion of an interaction.

Once a choice has been made, the two experiments are related by a restriction map. On the left of the map put smooth solutions of an independently specified whole problem. On the right put regional solutions satisfying transmission equations obtained from the regional actions. Restriction is useful when it goes both ways: every such pair comes from one whole solution, with the corresponding statement about changes of gauge. Requiring equality of \emph{all} derivatives on the right would often make the assertion immediate. The physical question is whether the much smaller amount of information exchanged by the equations is sufficient. Chapter~\ref{ch:07} proves that it is for a class of nonlinear equations and constructs a counterexample when a branch assumption is dropped.

A detector asks a different question of the map. It measures the change of a functional of the field when a source is switched on. Its answer involves the retarded inverse of the linearized equations, not merely the set of solutions of the homogeneous equations. A loop that crosses the cut asks a third question. Even if the two parties agree on every bulk equation and every response, what must each record at its two exposed endpoints so that the loop can be measured at all? These are reasons to keep \emph{fields}, \emph{causal response} and \emph{observable algebras} separate long enough to construct a map between them.

\section*{What a region carries}
A region carries a spacetime with its labeled faces, the bundles on which the fields live, fixed couplings and backgrounds, and a local action including the terms that specify the physical boundary problem. The word \emph{open} will mean that the data on an artificial cut may vary. It will never mean that boundary conditions at infinity, or at a real wall, have been discarded. A \emph{transparent seam} means that, after matching, there is no physical sheet left there. A source confined to a sheet or a field propagating along it is allowed, but it specifies a different target theory. A jump in a derivative can be precisely what that theory demands.

We use signature $(-,+,\ldots,+)$. A coorientation chooses which side of a hypersurface is positive. A collar is a thin strip of coordinates next to a face. A null face has no unit metric normal; none is needed to integrate an oriented flux. On a space of solutions $\delta$ denotes variation of the fields, rather than spacetime differentiation. The sign convention for the integrated presymplectic form is $\omega=\delta\theta$, where $\theta$ includes its corner contribution. The sign of a causal propagator belongs to its operator: throughout, the difference for a specified Hessian is its retarded inverse minus its advanced inverse. Multiplying the wave equation by $-1$ also multiplies that inverse by $-1$.

\section*{How to use the book}
If you know the action principle and wave equations from QFT and GR, you can start at Chapter~\ref{ch:01}. No previous use of jet bundles, gauge groupoids or microlocal analysis is assumed. For a short first reading, follow one scalar field through the signed collar in Chapter~\ref{ch:03}, the open variation in Chapter~\ref{ch:04}, and the transmission and response problems in Chapters~\ref{ch:07} and~\ref{ch:08}. Then return to the gauge theory chapters: what was an equality of field components becomes an equality after transport, and some transformations that look like redundancies on a closed slice carry a measurable charge at its exposed end. Exercises throughout have short checks or hints. The appendices develop three tools at the points where precision matters most: orientation, trace and support, and products of singular quantum fields.

The quantum question has a useful control experiment: sew the classical linearized theory first and quantize its observables. This produces a comparison algebra without multiplying two Hilbert spaces. A successful regional quantization should reproduce it within its stated scope. Reproducing its commutators does not choose a state; reproducing it at free order does not renormalize a pair of vertices that approach the seam from opposite sides. The last part starts with these distinctions. No Euclidean continuation is required.
